

\vspace{-0.05in}
\section{Conclusion}
\label{sec:conclusion}
\vspace{-0.05in}

We studied Bayesian optimisation with continuous approximations, by treating the
approximations as arising out of a continuous fidelity space.
While previous multi-fidelity literature has predominantly focused on a finite number of
approximations, \bocas applies to continuous fidelity spaces and can potentially
be extended to arbitrary spaces.
We bound the simple regret for \bocas and demonstrate that it is better than methods such
as \gpucbs which ignore the approximations and that the gains are determined by the
smoothness of the fidelity space.
When compared to existing multi-fidelity methods, \bocas is able to share information
across fidelities
effectively, has more natural modelling assumptions and has fewer hyper-parameters to tune.
Empirically, we demonstrate that \bocas is competitive with other baselines in synthetic
and real problems.

Going forward, we wish to extend our theoretical results to more general settings.
For instance, we believe a stronger bound on the regret might be possible
if $\phiz$ is a finite dimensional kernel.
Since finite dimensional kernels are typically not radial~\citep{sriperumbudur2016optimal},
our analysis techniques will not carry over straightforwardly.
Another line of work that we have alluded to is to study more general
fidelity spaces with an appropriately defined kernel $\phiz$.

